\section{Experiments}
\label{sec:experiments}

This section describes how the four families from Section~\ref{sec:methodology} are evaluated. It covers the recognition and interpretability metrics, the $63$-run experiment matrix that the report is built on, the ablation design that identifies which parts of the architecture and training procedure are load-bearing, and the attention-quality evaluation protocol used in Section~\ref{sec:results} to answer the central Darcet~\etal\ transferability question.

\subsection{Recognition Metrics}
\label{sec:metrics_reco}

Recognition quality is reported as \emph{character error rate} (CER) and \emph{word error rate} (WER). CER is the Levenshtein distance between the predicted string $\hat{y}$ and the reference $y$ normalized by the reference length:
\begin{equation}
\label{eq:cer}
\mathrm{CER}(\hat{y}, y) = \frac{d_{\mathrm{Lev}}(\hat{y}, y)}{|y|},
\end{equation}
computed per line and averaged over the split. WER is the same quantity computed at word granularity after whitespace tokenization. Both are reported at the \emph{best validation-CER epoch} of each run, following the convention of Retsinas~\etal~\cite{retsinas2022htrbestpractices}. Unless explicitly noted, ``test~CER'' means test-split CER at the best validation-CER epoch, and ``SWA CER'' means the CER of the SWA-averaged model measured after the BatchNorm recalibration pass described in Section~\ref{sec:training} --- which the trainer logs to \texttt{training.log} rather than to the per-epoch \texttt{results.csv}.

CER and WER alone do not fully answer the register-token question. They are used in Section~\ref{sec:results} to characterize (a) the near-flat CER-versus-$R$ curve and (b) the seed-noise floor against which any register-count effect must be judged. The interpretability side of the question is evaluated separately (Section~\ref{sec:attn_eval}).

\subsection{Experiment Matrix}
\label{sec:matrix}

The report is built on $63$ training runs, numbered \texttt{run\_143} through \texttt{run\_205}. Table~\ref{tab:matrix} groups them into six experiment blocks. IAM and READ2016 runs use $150$ epochs and $\mathrm{seed} = 42$ unless otherwise noted; synthetic runs targeted $15$ epochs but were truncated by the $24$-hour SLURM walltime to $2$--$11$ epochs depending on architecture cost, and are disclosed as such throughout.

\begin{table}[t]
\centering
\small
\caption{Experiment matrix. Runs marked $\dagger$ were truncated by the $24$\,h SLURM walltime; all other runs completed the full $150$-epoch budget.}
\label{tab:matrix}
\resizebox{\linewidth}{!}{%
\begin{tabular}{llcp{4.4cm}}
\toprule
Block & Runs & \# & Purpose \\
\midrule
Register sweep (IAM) & 143--156, 172 & 12 & CNN--RNN baseline + ViT-RGTS~v2 at $R \in \{0,2,4,8,16\}$, with and without SWA \\
Pretrained lanes (IAM) & 153--155, 173--174, 177--178, 180, 182--183 & 10 & ViT-B/16 (naive / two-group LR / frozen / gradual unfreeze); TrOCR (frozen / LLRD / gradual unfreeze / +SWA) \\
Architecture ablation & 157--165, 175, 179 & 12 & CNN stem, dual head, RNN depth, embedding dim, encoder depth, dropout, augmentation profile, legacy v1 layout \\
Training strategy ablation & 167--171, 176, 181, 183, 186, 188 & 11 & Base LR, batch/accumulation, SWA start epoch, SWA LR, dropout; seeds $\{42, 123, 456\}$ for reproducibility \\
READ2016 generalization & 166, 184--185, 187, 190 & 5 & CNN--RNN + ViT-RGTS~v2 at $R \in \{0,4,8,16\}$; HTR-VT run \texttt{run\_142} included as external comparator \\
Synthetic pretraining$^\dagger$ & 189, 191--205 & 16 & CNN--RNN + ViT-RGTS~v2 + ViT-B/16 + TrOCR on ${\sim}950$K synthetic lines, evaluated on IAM test \\
\midrule
\textbf{Total} & 143--205 & \textbf{63} & \\
\bottomrule
\end{tabular}%
}
\end{table}

\subsubsection{Register sweep}

The primary experiment sweeps $R \in \{0, 2, 4, 8, 16\}$ on ViT-RGTS~v2, and runs each with and without SWA. The $R = 0$ configuration keeps every other architectural choice (CNN stem, depth $6$, dim $256$, dual head) constant, and is the reference against which register-added variants are compared. The CNN--RNN baseline is included in the same block with and without SWA so that the register-vs.-CNN--RNN comparison is at equal training budget and identical evaluation protocol.

\subsubsection{Pretrained lanes}

Both pretrained models are exercised across the four fine-tuning strategies of Section~\ref{sec:finetune} to separate the pretrained-representation quality question from the fine-tuning strategy question. The block includes a naive uniform-LR fine-tune ($\eta = 10^{-3}$, no group differentiation) as a deliberate failure case: it collapses catastrophically, and its collapse is the cleanest available demonstration of why the differential-LR / LLRD strategies exist.

\subsubsection{Architecture ablation}

The architecture ablation isolates one design decision at a time relative to the ViT-RGTS~v2 baseline (\texttt{run\_146}: $R = 4$, no SWA). Ablated variables are: the CNN stem (with vs.\ without --- the latter falls back to raw $16 \times 16$ patch embedding, row-major flattened); head type (dual vs.\ RNN-only vs.\ CNN-only); RNN type (LSTM vs.\ GRU); RNN depth ($1$ vs.\ $3$ layers); embedding dim ($128$ vs.\ $256$ vs.\ $512$); encoder depth ($4$ vs.\ $6$ vs.\ $8$); dropout ($0.1$ vs.\ $0.3$); augmentation profile (moderate ``cnn'' vs.\ aggressive ``vit''); and the legacy v1 architecture with row-major patch flattening. The v1 layout is included precisely because it independently reproduces the CNN-stem-removal failure mode --- two convergent failures on the same underlying cause carry more weight than one.

\subsubsection{Training-strategy ablation}

The training-strategy ablation varies base LR ($5 \times 10^{-4}$ vs.\ $10^{-3}$), physical batch size and gradient accumulation ($8/4$ vs.\ $16/1$, both effective batch $\leq 32$), SWA start epoch ($75$, $110$, $135$ of $150$), SWA learning rate ($10^{-4}$ vs.\ $5 \times 10^{-4}$ vs.\ $10^{-3}$), and dropout ($0.1$ vs.\ $0.3$). Three seeds ($42$, $123$, $456$) on the $R = 4$ no-SWA configuration establish the reproducibility floor used to interpret all other deltas.

\subsubsection{READ2016 generalization}

The READ2016 block re-runs the register sweep on a different script (historical German), a different vocabulary size ($88$ vs.\ $79$ classes), and a larger training set ($8{,}349$ vs.\ $6{,}482$ lines). Its purpose is to test whether the register-count-vs-CER pattern observed on IAM is an IAM-specific artifact or a property of the architecture and loss. HTR-VT (\texttt{run\_142}) is included in the same block as an external comparator with an independently designed CTC ViT.

\subsubsection{Synthetic pretraining}

The synthetic block trains each of the four families on the ${\sim}950$K-line synthetic corpus (Section~\ref{sec:datasets}) and evaluates on real IAM test without further IAM fine-tuning. Its purpose is to quantify the render-to-real domain gap for this specific synthetic-data pipeline. All runs in this block hit the $24$\,h SLURM walltime before reaching the target $15$ epochs; the reported numbers are lower bounds on convergence, not converged results.

\subsection{Ablation Design Rationale}
\label{sec:abl_rationale}

Every ablation follows the same discipline: change exactly one variable from the ViT-RGTS~v2 baseline (\texttt{run\_146}), keep every other seed and configuration identical, train for $150$ epochs (or until the walltime for synthetic runs), and evaluate on the same IAM test split with the same greedy CTC decoder. When an ablation flips more than one variable at once (\eg\ ``ViT augmentation'' also changes the augmentation library's noise probability along with the geometric-transform strength), the additional variables are noted in Section~\ref{sec:results} where the row is presented, so the reader can see what is actually varying.

The report uses two independent significance criteria to interpret ablation deltas: the seed-noise floor of $\pm 0.08$~pp CER established by the three-seed reproducibility experiment, and a pairwise bootstrap of per-line CER differences against the baseline where the ablation-delta is comparable in magnitude to the seed floor. Any effect smaller than roughly $2 \times$ the seed-noise floor (\ie\ $\lesssim 0.16$~pp) is treated as unresolvable at this experimental budget and reported as such rather than being over-interpreted.

\subsection{Attention-Quality Evaluation}
\label{sec:attn_eval}

The register-token intervention is designed to change what attention maps look like more than what CER measures. Three complementary attention-side evaluations are used in Section~\ref{sec:results}, all applied to a fixed sample of $8$ IAM test lines (chosen for their diverse length, writer style, and character content, and re-used across all register counts to keep the visual comparison fair):

\begin{enumerate}
    \item \textbf{Per-character attention maps} in the style of Nikolaidou~\etal\ Fig.~5~\cite{nikolaidou2024beyondmem}. For each recognized character $c$ at CTC decoding column $t_c$, the attention row from the corresponding patch token in the final Transformer layer is reshaped to the $(H_p, W_p)$ grid, upsampled to the input resolution, and overlaid on the input line with the \texttt{inferno} colormap at $\alpha = 0.5$. Because our line images have a $1{:}8$ aspect ratio rather than the near-square word crops used by \emph{Beyond Memorization}, the raw overlay is complemented by a per-character crop of $\pm 12$ grid columns around $t_c$, with white patch-grid lines and a cyan CTC-column marker overlaid for reference.
    \item \textbf{Register absorption analysis.} For each layer $\ell$ we compute the fraction of total attention mass directed at register tokens versus patch tokens, and the L2 norm of each token position in the final-layer output. This yields two quantities per register count: the register share of attention $\alpha_{\mathrm{reg}}^{(\ell)} = \sum_{s \leq R} \|\mathbf{A}^{(\ell)}[:,:,:,s]\|_1 / \|\mathbf{A}^{(\ell)}\|_1$ (\ie\ what fraction of queries attend to registers), and the concentration index of the top-$k$ patch tokens by norm (\ie\ whether a small number of ``sink'' patches still exist when $R > 0$). These metrics operationalize the claim of Darcet~\etal\ in a way that CER cannot.
    \item \textbf{CTC-column concentration.} For each recognized character we measure how close the attention peak sits to the CTC decoding column: $\Delta_c = |\arg\max_{s} \mathbf{A}^{(L)}_{t_c, s} - t_c|$ in patch-column units, and the fraction of characters with $\Delta_c \leq 2$. This is the closest quantitative analogue to a ``localization accuracy'' for CTC self-attention and gives a single number that is monotonically better for cleaner attention.
\end{enumerate}

The third metric depends on a raw-attention anchoring correction described in Section~\ref{sec:results}: the raw self-attention of the CTC query column is contaminated by fixed high-norm sink patch columns, so a Gaussian column-prior centered at $t_c$ is used to suppress these when reporting $\Delta_c$. This correction is applied uniformly across all register counts so that it does not bias the comparison, and its magnitude is itself a diagnostic --- if registers were fully absorbing the sinks, the correction would be unnecessary.
